\documentclass[10pt]{article}
\usepackage[margin=1in]{geometry}
\usepackage{setspace}
\usepackage{hyperref}
\hypersetup{colorlinks=true, allcolors=black}
\setlength{\parindent}{0pt}
\setlength{\parskip}{0.55em}
\onehalfspacing
\pagestyle{empty}

\begin{document}

\begin{center}
HSU Ke Cheng\\
PhD Candidate in Economics, Hong Kong University of Science and Technology\\
Research Statement
\end{center}

My research interests lie in digital economics and industrial organization. My current and future work centers on one broad theme: how a decision made before uncertainty is resolved interacts with the adjustment that follows. An agent must commit before the uncertainty is resolved, and adjusting afterward is costly. Anticipating that cost changes the earlier decision, so the two are determined together. The theme splits into two streams. In one, consumers face uncertainty about match value. In the other, sellers face uncertainty about demand.

Consumers can resolve the match before paying, by searching, or after paying, by returning. Between search and purchase, product design shapes the distribution of match value and directs which consumers search which products. I study that margin with Professor AU Pak Hung. Between purchase and return, the buyer has already paid, and a return policy sets the cost of returning. The return rate feeds back into the price, which ties purchase efficiency to return efficiency. With Ce Bian and Liang Zhao, I follow the same feedback into pass-through and incidence. A calculation that keeps only the purchase margin can reverse the welfare conclusion.

For sellers, inventory is chosen before demand is realized. If demand exceeds inventory, the product stocks out. If inventory exceeds demand, the seller returns the unsold units. With Professor Kohei Kawaguchi and Professor Kei Kawai, I study resale price maintenance in Japanese publishing. The retail price is set in advance, so unsold copies are returned upstream and discarded, at an environmental cost.

The paragraphs below describe these papers, and then the questions I plan to take up next.

\textbf{Job Market Paper}

In my job market paper, ``Designing Return Policy in E-Commerce: A Trade-off between Purchase and Return Efficiency,'' I show that a return policy cannot be judged by whether consumers return efficiently ex-post, but by how purchase and return efficiency interact. Because the buyer does not bear the seller's return-handling cost, returns can exceed the efficient level. Tightening the return policy curbs those excessive returns, but also makes a higher price more attractive, which can exclude purchases that would have generated positive social surplus.

Specifically, I study return-shipping insurance on Chinese e-commerce platforms by estimating a structural model of purchase, insurance choice, and return under monopoly pricing across six apparel categories. Automatic coverage from the seller performs better when margins are high and handling costs are low, by expanding valuable purchases. Optional coverage from the consumer performs better when margins are low and handling costs are high, by disciplining returns. A menu that pairs a discounted price with a restocking fee balances the two margins and outperforms both for most categories.

\textbf{Other Papers}

In ``Pass-Through and Incidence with Product Returns,'' we study pass-through and incidence with product returns. In contrast to the case without returns, a tax such as a platform commission fee or a sales tax can create negative pass-through and improve social welfare. When a high price significantly raises consumers' return probability, the seller reduces the quantity supplied because of a lower effective margin. This downward-sloping supply curve can lower the equilibrium price in response to a tax increase, raising social surplus by reducing the return rate.

In ``Product Design with Heterogeneous Search Costs,'' we build on Wolinsky's (1986) sequential search model and introduce search-cost heterogeneity. Search costs shape not only whether consumers search but which products they search. In equilibrium, low-search-cost consumers seek niche products and high-search-cost consumers seek standardized products. The main result is that a fall in search costs can raise both prices when the endogenous shift in which product to search dominates the search-encouraging effect.

In ``Too Much Unsold Stock?: Environmental Consequences of Price Commitment,'' we study resale price maintenance in the Japanese publishing industry, where publishers set retail prices and unsold books are returned upstream and discarded. Resale price maintenance is procompetitive because it encourages sellers to carry inventory in the low-demand state. When realized demand falls short of that inventory, however, the seller cannot lower the price to clear the market, so unsold units pile up at an environmental cost. We build a dynamic oligopoly model with demand uncertainty, estimate it on monthly bookstore data, and ask how much unsold stock is due to price commitment and how a deposit-refund tax can reduce waste.

\textbf{Future Research Agenda}

My future research agenda spans three directions rooted in designing return policy.

\textbf{Search, purchase, and return jointly.} I plan to study search, purchase, and return jointly. My two consumer-side papers treat search and return in isolation, but search resolves match value only partially, so the two margins interact. In particular, I ask how search technology such as personal recommendation changes consumers' return behavior, and how the return policy in turn changes how consumers search.

\textbf{Bracketing, substitutes, and complements.} Allowing consumers to buy products of different sizes or colors creates a bracketing issue: a consumer orders several variants, keeps only one, and returns the rest, generating a large return loss. Bracketing is a special case of substitutes. I plan to study whether a policy that encourages the purchase of complements can solve the issue. Encouraging complements makes demand less elastic and strengthens the seller's ability to extract surplus. At the same time, returning one product makes the consumer more likely to return the other, so the seller can offer higher quality or a lower price to deter returns, partially offsetting that incentive.

\textbf{Return policy on a two-sided platform.} A return policy changes which consumers join the platform and how each of them returns, which changes how much a seller can earn and in turn selects which sellers participate and what prices they set. I plan to study how the composition on one side feeds into the other, and how return policy should be designed taking the cross-side network effect into account. In particular, should different instruments be delegated to different parties? Currently the refund rule is largely set by the platform, while discounts and prices are chosen by the seller. I ask whether the refund instrument should also be given to the seller.

\end{document}
